\unit{One atom, one photon: Jaynes--Cummings physics}
\unitref{W6L1; W6L2; Week 6 tutorial}
\subsection{Why a two-level atom is a useful model}
An atom has many levels, but when a narrow-band
cavity is tuned close to a single transition
\(\ket g\leftrightarrow\ket e\), the other
levels may be neglected to a good approximation.
The atom then behaves like a spin-\(1/2\)
system with
\[
\sigma_+=\ket e\bra g,\quad
\sigma_-=\ket g\bra e,\quad
\sigma_z=\ket e\bra e-\ket g\bra g.
\]
One photon of frequency \(\omega_c\) has
the right energy to excite the atom if
\(\omega_c\approx\omega_a\).
The dipole interaction
\(H_{\rm int}=-\mathbf d\cdot\mathbf E\)
contains products of photon and atomic
raising/lowering operators.

\subsection{Rotating-wave approximation}
Terms \(a\sigma_+\) and \(a^\dagger\sigma_-\)
exchange excitation between atom and field.
The counter-rotating combinations
\(a\sigma_-\) and \(a^\dagger\sigma_+\)
change excitation number by two and
oscillate rapidly near resonance.
If \(g\ll\omega_a,\omega_c\), the
rotating-wave approximation retains
the slowly varying terms:
\[
\boxed{H_{\rm JC}=
\hbar\omega_c a^\dagger a+
\hbar\omega_a\sigma_+\sigma_-+
\hbar g(a\sigma_++a^\dagger\sigma_-).}
\]
The excitation operator
\(N=a^\dagger a+\sigma_+\sigma_-\)
commutes with \(H_{\rm JC}\):
excitations move, but their total number is conserved.

\subsection{Reduce the problem to a two-state block}
For a fixed total excitation,
the coupled states are
\(\ket{e,n}\) and \(\ket{g,n+1}\).
After subtracting a common energy, the
Hamiltonian block is
\[
\frac{H_n}{\hbar}=
\frac12\begin{pmatrix}
\Delta&2g\sqrt{n+1}\\
2g\sqrt{n+1}&-\Delta
\end{pmatrix},
\quad \Delta=\omega_a-\omega_c.
\]
Its eigenfrequency separation is
\[
\boxed{\Omega_n=
\sqrt{\Delta^2+4g^2(n+1)}}.
\]
The factor \(\sqrt{n+1}\) comes from
\(a^\dagger\ket n=\sqrt{n+1}\ket{n+1}\).
It is a directly observable signature
of a quantised cavity field.

\begin{worked}{Vacuum Rabi oscillation}
Begin in \(\ket{e,0}\) at resonance
(\(\Delta=0\)). The interaction block is
\(\hbar g\sigma_x\), so
\[
\ket{\psi(t)}
=\cos(gt)\ket{e,0}
-i\sin(gt)\ket{g,1}.
\]
Thus
\[
\boxed{P_e(t)=\cos^2(gt),\quad
P_{g,1}(t)=\sin^2(gt)}.
\]
At \(t=\pi/(2g)\) the atom has emitted
one photon into the cavity.
At a full period \(t=\pi/g\),
the excitation has returned to the atom.
\end{worked}

\subsection{From one atom to many: Tavis--Cummings}
For \(N\) identical emitters with the same
coupling phase,
\[
H_I=\hbar g(aJ_++a^\dagger J_-),\quad
J_\pm=\sum_{i=1}^N\sigma_\pm^{(i)}.
\]
For two atoms with one excitation, define
\[
\ket S=\frac{\ket{eg}+\ket{ge}}{\sqrt2},
\quad
\ket A=\frac{\ket{eg}-\ket{ge}}{\sqrt2}.
\]
Because
\[
J_-\ket S=\sqrt2\ket{gg},\qquad
J_-\ket A=0,
\]
the symmetric state is \emph{bright},
with coupling \(\sqrt2g\);
the antisymmetric state is \emph{dark}
to the common cavity mode.

\subsection{Dicke states and superradiant enhancement}
For \(N\) symmetric spins, collective
angular-momentum labels \(J=N/2,m\)
describe a Dicke ladder. The lowering operator
has matrix element
\[
J_-\ket{J,m}=
\sqrt{(J+m)(J-m+1)}\ket{J,m-1}.
\]
If all atoms couple to a \emph{common}
radiation reservoir under appropriate
conditions, these factors enhance
collective spontaneous emission rates.
The coherent cavity-exchange enhancement
\(\sqrt N\,g\) and cooperative radiative
superradiance are related but not
identical phenomena: the reservoir and
initial state also matter.

\begin{worked}{Two-atom conservation check}
Start with \(\ket{eg,0}\).
It equals \((\ket{S,0}+\ket{A,0})/\sqrt2\).
The bright half oscillates with
\(\sqrt2g\), while the dark half stays.
Writing \(\vartheta=\sqrt2gt\),
\[
\ket{\psi(t)}
=\frac{\cos\vartheta\ket{S,0}+\ket{A,0}
-i\sin\vartheta\ket{gg,1}}{\sqrt2}.
\]
The mean cavity photon number is
\(\langle n\rangle=\tfrac12\sin^2\vartheta\).
The missing probability remains as atomic
excitation, preserving \(\langle N\rangle=1\).
\end{worked}

\begin{checkpoint}
To solve cavity-QED problems, identify
the conserved excitation number first.
That reduces an infinite-dimensional
Hamiltonian to independent \(2\times2\)
or small collective subspaces.
\end{checkpoint}
